% Teaser, Variant 2 — "break the loop" (cycle-centric). Column-width;
% \input inside a figure environment. The bias feedback loop drawn as a
% circuit (city -> data -> model -> dispatch -> city), with FAMAIL as the
% accent branch that exits the loop at the DATA node. Sibling of
% figures/figure-1: same palette + mark vocabulary, plain schematic grids,
% no mechanism. This figure is the WHY; fig:overview is the HOW.
%
% Only two numbers appear, both committed:
% % src: PAPER/supply-lift/data/a10/shz_a10_external_fairness.json
%        metrics.MigrantRatio.district_extremes.disparate_impact.before
%        = 0.3325 -> advantaged/disadvantaged service ratio 1/0.3325 = 3.0x
% % src: s10 corpus metrics.json — ~9,900 edited of ~100k trajectories (~10%);
%        intro/abstract: k=10,000, "about a tenth of the corpus"
\definecolor{figink}{HTML}{1A1A1A}
\definecolor{figgray}{HTML}{6F6F6F}
\definecolor{figfaint}{HTML}{C9C9C9}
\definecolor{figaccent}{HTML}{0F62FE}
\definecolor{figtrim}{HTML}{D55E00}
\begin{tikzpicture}[x=1cm, y=1cm, line cap=round, line join=round,
  pick/.pic={\draw[line width=0.8pt] (-0.09,-0.09) -- (0.09,0.09)
             (-0.09,0.09) -- (0.09,-0.09);},
  taxi/.pic={\draw[figink, fill=white, line width=0.7pt]
             (-0.09,-0.09) rectangle (0.09,0.09);},
  taxinew/.pic={\draw[figaccent, fill=white, line width=0.9pt]
                (-0.09,-0.09) rectangle (0.09,0.09);
                \draw[figaccent, line width=0.8pt]
                (0.13,0.13) -- (0.23,0.13) (0.18,0.08) -- (0.18,0.18);},
  boundary/.style={figink!65, line width=0.9pt, dash pattern=on 3pt off 1.6pt},
  panelframe/.style={figgray, line width=0.6pt},
  citygrid/.style={figfaint, line width=0.35pt},
  flow/.style={-{Stealth[length=5pt, width=4pt]}, figink, line width=0.9pt},
  flowacc/.style={-{Stealth[length=5pt, width=4pt]}, figaccent,
                  line width=0.9pt},
  lbl/.style={fill=white, fill opacity=0.85, text opacity=1,
              inner sep=1.2pt, rounded corners=1pt}]

  % =========================================================================
  % the loop, clockwise: city -> data -> model -> dispatch -> city
  % =========================================================================
  % station: the divided city (top-left)
  \node[font=\tiny, text=figink] at (1.6,-0.18)
    {a divided city --- $3.0\times$ service gap};
  \draw[citygrid, xstep=0.4, ystep=0.4] (0.6,-1.6) grid (2.6,-0.4);
  \fill[black, opacity=0.08] (1.8,-1.6) rectangle (2.6,-0.4);
  \draw[boundary] (1.8,-0.4) -- (1.8,-1.6);
  \draw[panelframe] (0.6,-1.6) rectangle (2.6,-0.4);
  \foreach \x/\y in {0.8/-0.6, 1.2/-0.6, 0.8/-1.0, 1.6/-1.0}
    \pic at (\x,\y) {taxi};
  \pic[figgray] at (1.2,-1.4) {pick};
  \foreach \x/\y in {2.0/-0.6, 2.4/-1.0, 2.0/-1.4}
    \pic[figgray] at (\x,\y) {pick};

  % station: the recorded data (top-right)
  \draw[panelframe, rounded corners=2pt] (3.9,-1.9) rectangle (5.35,-0.45);
  \foreach \yy/\sgn in {-0.62/1, -0.82/-1, -1.02/1, -1.22/-1}
    \draw[figink!70, line width=0.5pt] (4.0,\yy)
      .. controls (4.42,\yy+\sgn*0.05) and (4.84,\yy-\sgn*0.05) .. (5.25,\yy);
  \fill[black, opacity=0.08] (3.96,-1.52) rectangle (5.29,-1.3);
  \draw[figink, line width=0.7pt] (4.0,-1.41)
    .. controls (4.42,-1.36) and (4.84,-1.46) .. (5.25,-1.41);
  \node[lbl, anchor=east, font=\tiny, text=figink] at (5.27,-1.41)
    {${\sim}10\%$};
  \node[font=\tiny, text=figgray, align=center] at (4.62,-1.68)
    {driver\\trajectories};

  % station: the imitation model (bottom-right of the loop)
  \draw[panelframe, rounded corners=2pt] (3.9,-4.0) rectangle (5.35,-3.25);
  \node[font=\tiny, text=figink, align=center] at (4.62,-3.62)
    {imitation\\model};

  % station: dispatch (bottom-left)
  \draw[panelframe, rounded corners=2pt] (0.6,-4.0) rectangle (2.4,-3.25);
  \node[font=\tiny, text=figink, align=center] at (1.5,-3.62)
    {dispatch guided\\by the policy};

  % circuit arrows
  \draw[flow] (2.7,-1.0) -- (3.8,-1.0);
  \node[anchor=south, font=\tiny, text=figgray] at (3.25,-0.95) {recorded};
  \draw[flow] (4.62,-1.98) -- (4.62,-3.17);
  \node[anchor=west, font=\tiny, text=figgray] at (4.74,-2.58) {imitated};
  \draw[flow] (3.8,-3.62) -- (2.5,-3.62);
  \node[anchor=north, font=\tiny, text=figgray] at (3.15,-3.72) {deployed};
  \draw[flow] (1.5,-3.17) -- (1.5,-1.68);
  \node[anchor=east, font=\tiny, text=figgray, align=right] at (1.42,-2.5)
    {supply sent\\where it went};

  % the compounding, named at the loop's center
  \draw[-{Stealth[length=3.5pt, width=3pt]}, figgray, line width=0.7pt]
    (3.15,-2.28) ++(40:0.2) arc (40:-260:0.2);
  \node[font=\tiny, text=figgray, align=center] at (3.15,-2.75)
    {bias\\compounds};

  % =========================================================================
  % the FAMAIL branch: exit the loop at the data
  % =========================================================================
  \draw[flowacc] (5.35,-1.0) -- (5.75,-1.0);
  \draw[figaccent, line width=0.9pt] (5.75,-1.6) rectangle (8.35,-0.4);
  \node[font=\tiny, text=figink, align=center] at (7.05,-1.0)
    {\textbf{FAMAIL: edit the ${\sim}10\%$}\\
     \textbf{where the bias lives}\\
     kept $100\%$ real, generated $0$};
  \draw[flowacc] (7.05,-1.6) -- (7.05,-2.4);
  \draw[panelframe, rounded corners=2pt] (6.15,-3.15) rectangle (7.95,-2.4);
  \node[font=\tiny, text=figink, align=center] at (7.05,-2.78)
    {upweighted\\imitation};
  \draw[flowacc] (7.05,-3.15) -- (7.05,-3.47);

  % the exit: a city the policy serves on both sides
  \draw[citygrid, xstep=0.4, ystep=0.4] (6.05,-4.75) grid (8.05,-3.55);
  \fill[black, opacity=0.08] (7.25,-4.75) rectangle (8.05,-3.55);
  \draw[boundary] (7.25,-3.55) -- (7.25,-4.75);
  \draw[panelframe] (6.05,-4.75) rectangle (8.05,-3.55);
  \foreach \x/\y in {6.25/-3.75, 6.65/-4.15, 6.25/-4.55}
    \pic at (\x,\y) {taxi};
  \pic[figgray] at (7.05,-4.15) {pick};
  \foreach \x/\y in {7.45/-3.75, 7.85/-4.55}
    \pic[figgray] at (\x,\y) {pick};
  \pic at (7.45,-4.15) {taxinew};
  \pic at (7.85,-3.75) {taxinew};
  \node[font=\tiny, text=figink] at (7.05,-5.0)
    {policy serves both districts};
\end{tikzpicture}
